\chapter{Macular Pucker}
\index{Macular Pucker}
\label{sec:Macular Pucker}

\section{Overview}
\begin{itemize}[noitemsep]
  \item Macular pucker, also called epiretinal membrane or cellophane maculopathy, is a thin scar-like layer that forms on the surface of the retina over the macula
  \item It is a common condition in people over 50, affecting about 5--10\% of older adults
  \item It causes mild-to-moderate visual distortion (metamorphopsia) and blurring but rarely causes severe vision loss
\end{itemize}
\section{Causes}
This condition happens because glial cells migrate onto the retinal surface and form a contractile membrane over the macula. The membrane contracts, wrinkling and distorting the underlying retina. It can be primary (idiopathic, age-related) or secondary to posterior vitreous detachment, retinal tears or detachment, inflammation, trauma, or eye surgery (especially cataract and retinal surgery).
\section{Risk Factors}
\begin{itemize}[noitemsep]
  \item Age over 50 years
  \item Prior posterior vitreous detachment
  \item Eye surgery, especially retinal detachment repair
  \item Eye trauma
  \item Intraocular inflammation (uveitis)
  \item Diabetes and vascular retinal disease
\end{itemize}
\section{Signs and Symptoms}
\subsection{Common symptoms}
\begin{itemize}[noitemsep]
  \item Mild blurred or distorted central vision (straight lines appear wavy)
  \item Slight difficulty reading fine print
  \item Mildly reduced visual acuity
  \item Most cases are mild and many people have no noticeable symptoms
\end{itemize}
\subsection{Emergency warning signs}
\begin{itemize}[noitemsep]
  \item Sudden worsening of vision, new flashes of light, or a shadow or curtain in the vision
\end{itemize}
\section{Progression and Stages}
Macular pucker develops slowly over months to years. Most cases remain mild and stable, requiring no treatment. The membrane may contract progressively in some people, causing increasing distortion; severe cases with significant metamorphopsia or reduced acuity are the ones considered for surgery. Spontaneous improvement is rare but can occur.
\section{Complications}
\begin{itemize}[noitemsep]
  \item Progressive distortion and blurring of central vision
  \item Macular edema beneath the membrane
  \item Reduced quality of life from visual distortion
  \item Rarely, the membrane may cause a macular hole or tractional retinal detachment
\end{itemize}
\section{Diagnosis}
\begin{itemize}[noitemsep]
  \item Dilated eye examination
  \item Optical coherence tomography (OCT) is the key test, showing the membrane and retinal wrinkling
  \item Fluorescein angiography if vascular leakage or macular edema is suspected
  \item Visual acuity testing and Amsler grid testing to document distortion
\end{itemize}
\section{Prevention}
\begin{itemize}[noitemsep]
  \item Macular pucker is usually not preventable
  \item Managing underlying conditions such as diabetes and inflammation
  \item Prompt treatment of retinal tears and detachments
  \item Routine eye examinations to detect early changes
\end{itemize}
\section{Treatment Options}
\begin{itemize}[noitemsep]
  \item Observation for mild cases with regular monitoring
  \item Vitrectomy with membrane peeling for symptomatic cases with significant distortion or reduced vision
  \item Treatment of any associated macular edema (e.g., anti-VEGF or steroid therapy)
  \item Low-vision aids for persistent visual symptoms
  \item Most people do not require surgery; the decision is based on functional impairment
\end{itemize}
\section{Living with It}
\begin{itemize}[noitemsep]
  \item Use good lighting and magnification aids for reading
  \item Follow your eye specialist's monitoring schedule
  \item Report any sudden visual changes promptly
  \item Seek support from low-vision services if daily activities are affected
\end{itemize}
\section{Outlook}
The outlook is generally good. Most macular puckers remain mild and stable without treatment. When surgery is performed, a majority of patients experience improved distortion and vision, though the membrane can occasionally recur.
